% Einstiegsblatt, erzeugt aus einstieg/erzeugen.py – setzen mit ../bauen.sh
\documentclass[a4paper, 9.5pt, parskip=half]{scrartcl}
\usepackage[ngerman]{babel}
\usepackage[margin=16mm, top=14mm, bottom=14mm]{geometry}
\usepackage{wtwinci}
\usepackage{csquotes}
\usepackage{multicol}
\graphicspath{{/home/thobug/AI-Projects/app4webtrees/docs/screenshots/}{/home/thobug/AI-Projects/app4webtrees/docs/}}
\pagestyle{empty}
\setlength{\columnsep}{8mm}
\begin{document}
\noindent\begin{minipage}[c]{0.13\linewidth}\includegraphics[width=\linewidth]{icon/icon-512.png}\end{minipage}\hfill
\begin{minipage}[c]{0.84\linewidth}\sffamily{\fontsize{26}{30}\selectfont\bfseries\color{wtPetrol}wtWin in fünf Minuten\par}\vspace{3pt}{\small\color{wtGrau}Dein webtrees-Stammbaum als Desktop-Programm · wtWin (Windows) · wtTux (Linux) · wtMac (macOS, zum Testen)\par}\end{minipage}
\vspace{4mm}\noindent{\color{wtPetrol}\rule{\linewidth}{0.8pt}}\vspace{2mm}
\begin{multicols}{2}
\section*{1 · Installieren}
\textbf{Windows:} \texttt{wtWin-….exe} von \href{https://wtwin.bgg-home.de}{wtwin.bgg-home.de} laden und starten. Der Installer fragt zuerst nach der Sprache, braucht keine Adminrechte und ersetzt eine ältere Fassung von selbst. Beim ersten Start „Weitere Informationen“ › „Trotzdem ausführen“ (nicht signiert).

\textbf{Linux:} \texttt{sudo apt install ./wttux\_…\_amd64.deb}

\textbf{macOS:} \texttt{…-arm64.dmg} (Apple-Chip) oder \texttt{…-x64.dmg} (Intel); beim ersten Start unter Systemeinstellungen › Datenschutz „Dennoch öffnen“.
\section*{2 · Verbinden oder allein arbeiten}
\textbf{Mit deinem webtrees:} Im Browser bei webtrees anmelden, Menü \menue{App} › \menue{Mit wtWin verbinden}. Das Programm übernimmt Adresse und Konto, fragt einmal nach und öffnet den Stammbaum. Voraussetzung: webtrees 2.2 mit dem Modul api4webtrees (in nas4webtrees schon dabei).

\textbf{Von Hand:} Adresse der webtrees-Seite eintragen, \menue{Verbinden}, mit Benutzername und Passwort anmelden. Es gelten dieselben Rechte wie auf der Website.

\textbf{Ohne Server:} Rechts auf dem Startbildschirm \menue{Stammbaum anlegen} oder \menue{Aus GEDCOM-Datei übernehmen …}. Danach arbeitest du ohne Server, Passwort und Internet; webtrees läuft unsichtbar im Hintergrund.
\section*{3 · Drei Aufbauten}
\textbf{Navigator} (\taste{Strg+1}): links die Kinder, in der Mitte die Person mit Ehen, rechts vier Generationen Vorfahren. Ein Klick macht jemand anderen zur Mitte.

\textbf{Baum im Mittelpunkt} (\taste{Strg+2}): Personenliste, Baum und Personentafel nebeneinander; Vollbild zum Arbeiten im Baum.

\textbf{Familienansicht} (\taste{Strg+3}): Vorfahren beider Partner oben, das Paar in der Mitte, Kinder mit Partnern und Enkeln unten.
\section*{4 · Erstellen und prüfen}
\textbf{Tafeln} (15 Arten, auch Plotter und Großdruck), \textbf{Listen} und \textbf{Bücher} (Ahnen, Nachfahren, Ortsfamilienbuch, auch als Website) im Menü \menue{Erstellen}, als PDF oder Druck.

\textbf{Extras:} Plausibilitätsprüfung mit 61 Regeln, Personen zusammenführen, Orte bereinigen, Datumsrechner, Kirchenjahr.

\textbf{Quellen, Orte, Archiv, Aufgaben:} eigene Fenster unter \menue{Ansicht} (\taste{Strg+5} bis \taste{Strg+9}).
\section*{5 · Hilfe}
\taste{F1} öffnet die Hilfe zum aktuellen Fenster. \taste{Strg+F} sucht eine Person, \taste{F5} lädt neu, \taste{Strg+C} kopiert den Personentext.

Handbuch als PDF, Fragen und Fehler: \href{https://wtwin.bgg-home.de}{wtwin.bgg-home.de} · \href{https://github.com/thobgg/app4webtrees/issues}{github.com/thobgg/app4webtrees/issues}
\vspace{3mm}
\noindent\bild{desktop-navigator.png}
\end{multicols}
\vspace{2mm}
\noindent{\color{wtLinie}\rule{\linewidth}{0.5pt}}\\[2pt]
{\sffamily\footnotesize\color{wtGrau}wtWin 1.52 · GPL-3.0 · wtwin.bgg-home.de}
\end{document}
